\scanpage{74}
\mspage{73}
\begin{flushright}31\end{flushright}
Soit $(Y,X)$ un $S$ -couple lisse de codimension $1$, et supposons que $Y$ soit défini dans $X$ par une équation $t=0$. Notons aussi par $t$ la section inversible de $\Gamma(U,\mathcal O_U)$ obtenu par restriction (notation de (3.1)). De la suite exacte de Kummer (IX.3.2) sur $U$
\diagram{\begin{tikzcd}0\arrow[r]&\mu_n\arrow[r]&\mathbb G_m\arrow[r,"n"]&\mathbb G_m\arrow[r]&0\end{tikzcd}}
on déduit un morphisme ``cobord''
\diagram{\begin{tikzcd}j_*\mathbb G_m\arrow[r,"\delta"]&R^1j_*\mu_n\end{tikzcd}}
d'où une section $\delta t$ de $R^1j_*\mu_n$. Puisque $n$ est annihilateur de $R^1j_*\mu_n$, cette section donne un morphisme
\[
\tag{3.9}
\varepsilon_t:(\mathbb Z/n)_Y\longrightarrow(R^2i^!)(\mathbb Z/n)_X\approx i^*(R^1j_*)(\mathbb Z/n)_U.
\]
Nous verrons que $\varepsilon_t$ est un bijection canonique.

\uline{Théorème(3.10):} Soit $C$ La famille des couples lisses de codimension constant. Il existe un fonction et un seul $\varphi$ donnant pour chaque $S$ -couple lisse $(Y,X)$ de codimension $c$ \inserted{et chaque $n$ premier aux caractéristiques résiduelles de $X$} un isomorphisme (dit canonique) de faisceaux
\deleted{\diagram{\begin{tikzcd}
(\unsure{Y|U},X)\ -\ :\ \mathbb Z/n\arrow[r]&\underline H_Y^{2c}(X,\unsure{\gamma|\mathbb Z}^{\unsure{u|n}}
\end{tikzcd}}}
